\part{Introduction et contexte général}
\chapter{Introduction générale}
\label{chap:introduction}

\subsection{La révolution silencieuse de l'apprentissage automatique}
L'aube du XXIe siècle est marquée par une transformation technologique aussi profonde que silencieuse : l'avènement de l'apprentissage automatique (AA) comme force motrice de l'innovation et de l'organisation sociale. Des diagnostics médicaux assistés par ordinateur à la gestion des chaînes logistiques mondiales, de la conduite de véhicules autonomes à la curation personnalisée de nos flux d'information, les algorithmes d'AA ne sont plus des curiosités de laboratoire. Ils forment une infrastructure invisible innovante qui médiatise de plus en plus nos échanges avec le monde, nos institutions et les autres.

Cette révolution est basée sur une notion apparemment simple : l'aptitude des machines à déceler des motifs à partir de données, sans nécessiter d'être spécifiquement codées pour chaque mission donnée. Cette compétence a ouvert la voie à des performances inédites en matière de prédiction et de prise de décision, notamment avec l'apparition des modèles d'apprentissage profond (Deep Learning). Toutefois, cette force a été obtenue à travers une complexité exponentielle. Les architectures de réseaux neuronaux profonds, qui possèdent souvent plusieurs milliards, voire des billions de paramètres (p. ex. GPT-5 \footnote{https://openai.com/fr-FR/research/}), fonctionnent fréquemment comme des systèmes en boîte noire : leurs mécanismes internes, qui transforment une entrée (par exemple, une image) en une sortie (par exemple, un diagnostic), sont si complexes qu'ils défient l'intuition et l'analyse humaine directe.

\subsection{La crise de la boîte noire : confiance, responsabilité et pouvoir}
Cette opacité n'est pas un simple défi technique ; elle engendre une véritable crise conceptuelle et sociétale. Comment pouvons-nous faire confiance à un système dont nous ne comprenons pas le raisonnement ? Comment attribuer la responsabilité en cas d'erreur ou de décision préjudiciable ? Comment s'assurer que ces systèmes, déployés à grande échelle, ne perpétuent pas ou n'amplifient pas les biais et les injustices existants dans nos sociétés ?

Ces questions sont au cœur de ce rapport. Elles nous obligent à dépasser une vision purement instrumentaliste de la technologie pour l'examiner comme un phénomène social et politique, porteur de nouvelles logiques de pouvoir et de contrôle. La «boîte noire» n'est pas seulement opaque, elle est puissante.\\ {\large Elle classifie}, {\large elle trie}, {\LARGE elle recommande}, {\huge elle décide}. Aussi, l'un des enjeux centraux de notre époque est qu'il devient « impossible pour un individu de mettre en doute une décision automatique, prise à l'encontre de ses intérêts, en demandant à l'algorithme de personnalisation de lui rendre des comptes sur les tenants et aboutissants de sa décision ».

Cette inaptitude de fournir des explications desdits systèmes, ce déficit de redevabilité/acceptabilité (\textit{accountability}), est le fil rouge qui relie les quatre grandes questions abordées dans ce rapport.

\subsection{Problématique et questions de recherche}
Cette crise de l'opacité n'est pas seulement technique ; elle est fondamentalement politique. Elle marque l'émergence d'une nouvelle forme de pouvoir que nous nommons le Léviathan \footnote{Inspiré par l'essaie philosophique de Thomas Hobbe de 1651 du même titre et soulignant la souveraineté absolue des états surpuissants, gouvernant de façon totalitaire leurs sujets. Ce livre, lui-même inspiré par le «monstre» légendaire de la bible hébraïque apparaissant dans l'ancien testament dans les livres de (Job 3:8, 40:25–41:26), Psaumes (Psaumes 74:14, 104:26), et Ésaïe (Ésaïe 27:1, 51:9)} algorithmique. La problématique centrale de cette recherche est que ce nouveau pouvoir fonctionne selon une logique qui est à la fois radicalement nouvelle dans sa forme technologique et profondément ancienne dans sa structure sociale : celle de la bureaucratie. Nous soutenons que les algorithmes de personnalisation constituent l'aboutissement d'une rationalité bureaucratique, telle qu'analysée par Max Weber et David Graeber, mais sous une forme privatisée, automatisée et opaque. C'est à travers ce cadre analytique de la bureaucratie algorithmique que nous aborderons les quatre questions structurantes qui forment les piliers de ce travail :

\begin{enumerate}[label=\textbf{Q\arabic*:}, wide, labelindent=0pt]
    \item \textbf{L'interprétabilité :} Comment définir, évaluer et mettre en œuvre l'interprétabilité des modèles d'AA, et pourquoi est-elle cruciale sur les plans éthique et pratique ?
    \vspace{2em}
    \item \textbf{Les systèmes de recommandation :} Quels sont les mécanismes qui sous-tendent ces outils omniprésents, et quels compromis leurs différentes architectures impliquent-elles ?
    \item \textbf{La bureaucratie algorithmique :} Quel lien peut-on établir entre la logique des algorithmes de personnalisation et le phénomène bureaucratique, tel qu'analysé par la sociologie classique et contemporaine ?
    \item \textbf{La responsabilité en pratique :} Comment passer de la théorie à la pratique en développant des cadres d'intervention pour modifier les algorithmes prédateurs et en formulant une déontologie pour leurs concepteurs ?
\end{enumerate}

\subsection{Structure du rapport}
Pour répondre à ces questions\footnote{Les questions originales, telles que reçues, sont à l'annexe F}, le rapport est structuré en plusieurs parties thématiques qui s'articulent de manière progressive.
\begin{itemize}[label=\S]
    \item La \textbf{partie II} établit les fondements théoriques en retraçant la trajectoire du concept de bureaucratie, de l'analyse de la rationalité chez Max Weber à la critique de la violence structurelle chez David Graeber, en passant par les apports de Béatrice Hibou. Cette revue de littérature approfondie construit le cadre sociologique essentiel pour l'ensemble de la thèse.
    \item La \textbf{partie III} opère la synthèse socio-technique centrale en développant le concept de «~bureaucratie algorithmique~». Elle met en dialogue les théories de la bureaucratie avec le fonctionnement des systèmes d'IA pour démontrer comment ces derniers intensifient et privatisent les logiques de pouvoir bureaucratiques.
    \item La \textbf{partie IV} passe de la critique à l'action en proposant des cadres concrets pour la responsabilité en pratique. Elle détaille une méthodologie d'intervention éthique pour modifier les algorithmes prédateurs et esquisse les principes d'un code de déontologie pour les concepteurs.
    \item Les \textbf{parties V et VI} plongent au cœur des fondements techniques qui sous-tendent l'analyse. Elles explorent en profondeur l'évolution des systèmes de recommandation (répondant ainsi à la Q2) et les enjeux de l'IA explicable (XAI, en lien avec la Q1), tout en illustrant les dangers de la bureaucratie algorithmique à travers des études de cas concrètes comme la \textit{Toeslagenaffaire} et la police prédictive.
    \item Enfin, la \textbf{partie VII} offre une conclusion générale, synthétisant les apports de la recherche, discutant de ses limitations et ouvrant des perspectives sur la nécessité d'une gouvernance démocratique du \textbf{«~Léviathan algorithmique~»}.
\end{itemize}

Ensemble, ces chapitres visent à offrir une analyse multifacettes, naviguant de la technique à la sociologie, de la théorie à la pratique, pour éclairer les enjeux de l'un des phénomènes les plus déterminants de notre temps.

\subsection{L'impact sociétal des systèmes algorithmiques}

L'intégration croissante des systèmes d'IA dans les rouages de la société ne se limite pas à des avancées techniques ; elle reconfigure profondément les structures de pouvoir, les modes de gouvernance, l'accès à l'information et le bien-être individuel. Cette section examine ces transformations à travers trois lentilles critiques : l'émergence d'une «bureaucratie algorithmique», les effets des écosystèmes de recommandation sur la diversité et le bien-être, et une étude de cas emblématique des préjudices algorithmiques.
